\emph{Preliminaries.} The modes are tower planes with $a^{*} = v$, $a = w$; by Theorem \claimref{tower.ladder-pairs-satisfy-car-vectors} and Fact \claimref{field-theory.imported-tower-limit} they satisfy the canonical anticommutation relations and act on the Fock space $\mathcal{F}$ with its positive inner product, for which $*$ is the adjoint. Everything below is finite-dimensional linear algebra on $\mathcal{F}$ with real coefficients; no unit is needed. Since the $Q_{i}$ are commuting self-adjoint operators, $\mathcal{F}$ is the orthogonal sum of the sectors $\mathcal{F}_{s}$, and since $H$ commutes with every $Q_{i}$ it maps each sector into itself, so $E(s)$, the lowest eigenvalue of $H$ on $\mathcal{F}_{s}$, is defined for every nonzero sector. If $X$ has charge $h$ then for $\psi \in \mathcal{F}_{s}$, $Q_{i}X\psi = XQ_{i}\psi + h_{i}X\psi = (s_{i} + h_{i})X\psi$, so $X\mathcal{F}_{s} \subset \mathcal{F}_{s+h}$; taking adjoints of $[Q_{i}, X] = h_{i}X$ gives $[Q_{i}, X^{*}] = -h_{i}X^{*}$, so $X^{*}$ has charge $-h$.

\emph{(i).} For $\psi \in \mathcal{F}_{s}$ and $h \ne 0$, $X\psi \in \mathcal{F}_{s+h}$, which is orthogonal to $\mathcal{F}_{s}$ (eigenspaces of a self-adjoint $Q_{i}$ with $h_{i} \ne 0$ for different eigenvalues), so $\langle\psi, X\psi\rangle = 0$. Each $Q_{i}$ commutes with $H$ and so maps each eigenspace of $H$ into itself; a one-dimensional eigenspace is then spanned by a common eigenvector of the $Q_{i}$, that is, it lies in a sector. More generally an eigenspace $V$ of $H$ is the orthogonal sum of its intersections $V_{s}$ with the sectors, and for $\psi = \sum_{s}\psi_{s}$ one has $\langle\psi, X\psi\rangle = \sum_{s}\langle\psi_{s+h}, X\psi_{s}\rangle$, which vanishes unless some $V_{s}$ and $V_{s+h}$ are both nonzero.

\emph{(ii).} Let $H\psi = E(s)\psi$, $\psi \in \mathcal{F}_{s}$, and $X\psi \ne 0$. Then $X\psi \in \mathcal{F}_{s+h}$ and $E(s+h)$ is at most its Rayleigh quotient, $E(s+h) \le \langle X\psi, HX\psi\rangle/\langle X\psi, X\psi\rangle$, with equality if and only if $X\psi$ is an eigenvector of $H$ for the eigenvalue $E(s+h)$ (finite dimension, $H$ self-adjoint on $\mathcal{F}_{s+h}$). Write $HX\psi = XH\psi + [H, X]\psi = E(s)X\psi + [H, X]\psi$; this gives the first bound. The second is the same argument for $X^{*}$, of charge $-h$.

\emph{(iii).} From $[H_{0}, X] = \varepsilon X$ with $\varepsilon$ real and $H_{0}$ self-adjoint, the adjoint gives $[H_{0}, X^{*}] = -\varepsilon X^{*}$. With $C = [X, X^{*}]$,
\[
[XX^{*}, X] = X[X^{*}, X] = -XC, \qquad [XX^{*}, X^{*}] = [X, X^{*}]X^{*} = CX^{*},
\]
so $[H, X] = \varepsilon X + GXC$ and $[H, X^{*}] = -\varepsilon X^{*} - GCX^{*}$. On $\psi \in \mathcal{F}_{s}$, $XC\psi = c(s)X\psi$, and $X^{*}\psi \in \mathcal{F}_{s-h}$ gives $CX^{*}\psi = c(s - h)X^{*}\psi$. Hence the two quotients of (ii) are exactly $\varepsilon + Gc(s)$ and $-\varepsilon - Gc(s - h)$, which are the two inequalities, and adding them bounds the second difference by $G(c(s) - c(s - h))$. If the ground state of $H$ lies in $\mathcal{F}_{s^{*}}$, then $E(s^{*}) \le E(s^{*} \pm h)$, so both differences are nonnegative, and their sum, hence each, is at most $G(c(s^{*}) - c(s^{*} - h))$.

\emph{(iv).} For $H_{4}$ take the charges $N_{a}$, $N_{c}$ and $X = P^{*} = \sum_{b}v_{c}v_{a}$, which raises each by one: $h = (1, 1)$. $H_{0} = \sum_{b}K(n_{a} + n_{c}) = KN$ with $N = N_{a} + N_{c}$ has $[N, P^{*}] = 2P^{*}$, so $\varepsilon = 2K$. By Proposition \claimref{field-theory.how-close-tower-comes-ccr}, $[P, P^{*}] = n_{b} - N$, so $C = [P^{*}, P] = N - n_{b}$, which on the sector $(N_{a}, N_{c})$ is the scalar $N_{a} + N_{c} - n_{b}$; on $(n, n)$, $c = 2n - n_{b}$, and $c(n, n) - c(n - 1, n - 1) = 2$. So (iii) applies with the bound $2G$. By Theorem \claimref{synthesis.tower-pairing-gap-at-fixed-charges-is-twice-the-quasiparticle-energy}(i), every level of the sector $(N_{a}, N_{c})$ is $K(2p + s_{a} + s_{c}) - G(j + j_{z})(j - j_{z} + 1)$ with $p$ pairs on $\Omega = n_{b} - s_{a} - s_{c}$ blocks; the last factor is largest at $j = \Omega/2$, where it is $p(\Omega - p + 1)$. So every level with $p$ pairs and $s = s_{a} + s_{c}$ singly occupied blocks is at least
\[
2Kp + Ks - Gp(n_{b} - s - p + 1) = E(p) + Ks + Gps \ge E(p), \qquad E(n) = 2Kn - Gn(n_{b} - n + 1),
\]
the band of Proposition \claimref{field-theory.composite-scalars-mean-field-level}(2), the lowest level of the sector $(n, n)$; its vector is $\psi_{n} = (P^{*})^{n}\Omega$, the quasi-spin state of Proposition \claimref{field-theory.how-close-tower-comes-ccr} with $j = n_{b}/2$, $j_{z} = n - n_{b}/2$. Hence the ground energy of $H_{4}$ on all of $\mathcal{F}$ is $\min_{n}E(n)$. Completing the square, $E(n) = Gn^{2} - (G(n_{b} + 1) - 2K)n = G(n - \gamma)^{2} - G\gamma^{2}$ with $\gamma = (n_{b} + 1)/2 - K/G$, so the minimum over integers $0 \le n \le n_{b}$ is at the integer $n^{*}$ nearest $\gamma$ (clamped to $0$ when $\gamma < \tfrac12$); since $\gamma \le (n_{b} + 1)/2 < n_{b} - \tfrac12$ for $n_{b} \ge 2$, $n^{*} \le n_{b} - 1$. The quasi-spin ladder gives $P^{*}\psi_{n} = \psi_{n+1}$, $P\psi_{n} = n(n_{b} - n + 1)\psi_{n-1}$, and $\|P^{*}\psi_{n}\|^{2} = (j - j_{z})(j + j_{z} + 1)\|\psi_{n}\|^{2} = (n + 1)(n_{b} - n)\|\psi_{n}\|^{2}$, $\|P\psi_{n}\|^{2} = (j + j_{z})(j - j_{z} + 1)\|\psi_{n}\|^{2} = n(n_{b} - n + 1)\|\psi_{n}\|^{2}$.

(a) $n^{*} \ge 1$ exactly when $\gamma > \tfrac12$, that is $Gn_{b} > 2K$ (at $Gn_{b} = 2K$, $E(0) = E(1)$). Then $E(n^{*} \pm 1) - E(n^{*}) = G((n^{*} \pm 1 - \gamma)^{2} - (n^{*} - \gamma)^{2}) = G(1 \pm 2(n^{*} - \gamma))$, and $|n^{*} - \gamma| \le \tfrac12$ puts both in $[0, 2G]$ with sum $2G$: the bound of (iii) holds with equality, as (ii) says it must, since $P^{*}\psi_{n^{*}}$ and $P\psi_{n^{*}}$ are lowest vectors of their sectors. $(P + P^{*})\psi_{n^{*}} = \psi_{n^{*}+1} + n^{*}(n_{b} - n^{*} + 1)\psi_{n^{*}-1}$, with the norms above. The pair-breaking threshold at fixed numbers is $Gn_{b}$ above the band (Theorem \claimref{synthesis.tower-pairing-gap-at-fixed-charges-is-twice-the-quasiparticle-energy}(ii)), so the larger rotor difference over the threshold is at most $2G/(Gn_{b}) = 2/n_{b}$; at the gap equation $Gn_{b} = 2E$ of (iii) there, the rotor is at most $4E/n_{b}$.

(b) If $Gn_{b} < 2K$ then $\gamma < \tfrac12$, $n^{*} = 0$, and every other level is positive: band levels $E(p) = p(2K - G(n_{b} - p + 1)) \ge p(2K - Gn_{b}) > 0$ for $p \ge 1$, and levels with $s \ge 1$ exceed $E(p) + Ks > 0$ ($K > Gn_{b}/2 > 0$). So $\Omega$ is the unique ground state, $P^{*}\Omega = \psi_{1}$ has energy $E(1) = 2K - Gn_{b} > 0$, and the sector $(1, 1)$ also contains the level $2K$ (one $a$ and one $c$ on different blocks, $p = 0$, $s_{a} = s_{c} = 1$). As $Gn_{b} \to 2K$ from below, $E(1) \to 0$.

(c) At $K/G = n_{b}/2 - n$ with $0 \le n \le n_{b} - 1$ an integer, $\gamma = n + \tfrac12$, so $E(n) = E(n + 1) = \min E$ and the ground eigenspace contains $\psi_{n}$ and $\psi_{n+1}$; for $\psi = \psi_{n} + \psi_{n+1}$, $\langle\psi, P\psi\rangle = \langle\psi_{n}, P\psi_{n+1}\rangle = \|P^{*}\psi_{n}\|^{2} = (n + 1)(n_{b} - n)\|\psi_{n}\|^{2} \ne 0$. The ground eigenspace is always the span of the band vectors $\psi_{n}$ with $E(n) = \min E$: a level with $s = 0$ and $j < n_{b}/2$ lies strictly above $E(p)$, since $(j + j_{z})(j - j_{z} + 1) = j(j + 1) - j_{z}(j_{z} - 1)$ increases strictly with $j$; a level with $s \ge 1$ is at least $E(p) + Ks + Gps$, which exceeds $\min E$ unless $K = 0$ and $p = 0$, and then it equals $E(0) = 0 > E(n^{*})$, because $K = 0$ gives $\gamma = (n_{b} + 1)/2 \ge \tfrac32$ and $E(n^{*}) = -Gn^{*}(n_{b} - n^{*} + 1) < 0$. So at every other coupling, where the minimiser $n^{*}$ is unique, the ground state is $\psi_{n^{*}}$ up to a factor, it lies in the sector $(n^{*}, n^{*})$, and by (i) it has $\langle P\rangle = 0$.

\emph{Reading (argued).} The uniform pair field is the helicity-zero composite scalar of the floor. Clauses (iv)(a) and (iv)(b) say exactly how its gap behaves: with a condensate ($n^{*} \ge 1$) its whole weight from the ground state sits on the two neighbouring band levels, within $2G$, and no weight at the threshold; this is the Goldstone rotor of the broken pair phase (the flat phase direction of Proposition \claimref{field-theory.composite-scalars-mean-field-level}(4) seen at finite size), and in the collective limit at fixed mean-field gap it closes as $4E/n_{b}$. Without a condensate the bound pair is gapped, at $2K - Gn_{b}$, and the gap closes where condensation begins. Clause (i) is what makes a finite order parameter need degeneracy across sectors, and (iii) is the finite form of the Goldstone theorem for a pairing-type element: it needs only $[H_{0}, X] = \varepsilon X$ and a commutator $[X, X^{*}]$ that is a function of the charges, so it covers several species paired by one uniform field with a common shell energy, and any extra term in $H_{0}$ that commutes with $X$. For a composite Higgs doublet, (i) applies to each charge the doublet carries; in a condensed doublet the directions of the broken charges are of this rotor kind, so a gapped composite doublet would have to be one that does not condense, or a non-uniform mode. The doublet floor and two-condensate couplings, where $[X, X^{*}]$ is no longer a function of the charges, are not treated here.
